\section{Extension E: Two valence electrons (calcium $4s4p$)}\label{sec:two-electron}
All earlier levels have one active electron. Alkaline-earth atoms such as calcium have two: an $s$ electron and a $p$ electron. They can combine into a triplet
($S=1$, levels $^3P_{0,1,2}$) or a singlet ($S=0$, level $^1P_1$). The code is \texttt{src/two\_electron.py}.

\subsection{Model}
The basis is $\ket{m_l}\otimes\ket{m_{s1}}\otimes\ket{m_{s2}}$ (12 states) and
\begin{equation}
  H = E_t + \Delta\,(\tfrac14 - \vb{s}_1\!\cdot\!\vb{s}_2) + \zeta\,\vb{l}\!\cdot\!\vb{s}_2 + \mu_B B\,(L_z + g_s(S_{1z}+S_{2z})).
\end{equation}
The second term is the exchange splitting: zero for the triplet and $\Delta$ for the singlet. The third is the spin--orbit coupling of the $p$ electron.
Within the triplet it gives the Land\'e interval rule, $E(^3P_2)-E(^3P_1) : E(^3P_1)-E(^3P_0) = 2:1$. It is \emph{not} diagonal in $S$,
so it mixes a little $^1P_1$ into $^3P_1$. The three parameters are fitted to the four measured term energies, which leaves one number as a consistency check.

\subsection{Results}
\begin{itemize}
  \item With the calcium energies from the NIST Handbook \cite{nistca} the fit reproduces all four levels to $0.09\ \si{cm^{-1}}$: the interval rule holds in calcium.
        The fitted values are $\zeta=105.3\ \si{cm^{-1}}$ and $\Delta=8388\ \si{cm^{-1}}$.
  \item The $g_J$ factors come out as $1.5$ for $^3P_1$ and $^3P_2$ and $1$ for $^1P_1$, read from the model's own eigenvalues at a tiny field.
  \item The $^1P_1$ admixture in $^3P_1$ is $\varepsilon^2 = 7.8\times10^{-5}$, equal to the perturbation-theory value $\bigl(\zeta/\sqrt2\,/\,\Delta E\bigr)^2$.
  \item That admixture lets $^3P_1$ decay to the ground state, the $657\ \si{nm}$ intercombination line. Scaling the $^1P_1$ decay rate by
        $\varepsilon^2$ and the cube of the wavelength ratio gives a lifetime of about $0.22\ \si{ms}$, against $0.39\ \si{ms}$ ($410\ \si{Hz}$ linewidth) and $0.42\ \si{ms}$ ($375\ \si{Hz}$) in two NIST-hosted calcium papers:
        the right order of magnitude and a factor $1.8$ to $1.9$ off. The estimate assumes one shared dipole matrix element and uses the NIST $^1P_1$ decay rate $2.20\times10^8\ \si{s^{-1}}$ \cite{nistca}, so it is a plausibility check, not a precise prediction.
\end{itemize}

\begin{figure}[h]
\centering
\includegraphics[width=\textwidth]{two_electron_calcium_levels.pdf}
\caption{Calcium $4s4p$ levels against field. Left: the three triplet levels. Right: Zeeman fans relative to each level's zero-field centre ($^3P_1$, $^3P_2$: $g\approx1.5$; $^1P_1$: $g=1$).}
\end{figure}

\subsection{Limits}
The $^3P_1$ energy comes from a measured frequency ($455\,986\,240\,494\,135.8\ \si{Hz}$ \cite{cafreq}, giving $15210.0638\ \si{cm^{-1}}$); the other three term energies and the $^1P_1$ decay rate are from the NIST Handbook tables \cite{nistca}, which give no uncertainties. Spin--spin interaction is not modelled
(it is large in helium, so this model suits calcium, not helium), nor is configuration mixing with $4s3d$ or hyperfine structure. Only levels are computed, not a line spectrum.
